\documentclass[conference]{IEEEtran}

\usepackage{cite}
\usepackage{amsmath,amssymb,amsfonts}
\usepackage{algorithmic}
\usepackage{graphicx}
\usepackage{textcomp}
\usepackage{xcolor}
\usepackage{booktabs}
\usepackage{tabularx}
\usepackage{tagging}
\usepackage{hyperref}

\usepackage{sc26repro}

\begin{document}

\twocolumn[%
{\begin{center}
\Huge
Appendix: Artifact Description/Artifact Evaluation
\end{center}}
]

%%%%%%%%%%%%%%%%%%%%%%%%%%%%%%%%%%%%%%%%%%%%%%%%%%%%%
%  AD Appendix
%%%%%%%%%%%%%%%%%%%%%%%%%%%%%%%%%%%%%%%%%%%%%%%%%%%%%

\appendixAD

\section{Overview of Contributions and Artifacts}

\subsection{Paper's Main Contributions}

\begin{description}
\item[$C_1$] A theoretical proof, via blocked CDAGs and the Red-Blue-White Pebble Game, that a poor data layout can inflate the I/O lower bound by up to a factor of the cache-line length $B$ relative to a layout-aware optimum, and that no schedule can close this gap. The full proof, cut from the paper for space, is bundled with the artifact at \texttt{Proofs/necessity\_proof.pdf}; it is not exercised by the software.
\item[$C_2$] A composable layout algebra of five primitives (pad, permute, block, shuffle, zip) whose compositions express any bijective array rearrangement.
\item[$C_3$] Two heuristic cost metrics: the average new-block count $\mu$ (guiding permutation) and the average block distance $\Delta$ (guiding blocking); plus a cookbook (Table~II) mapping observable performance symptoms to corrective primitives.
\item[$C_4$] Empirical validation on microbenchmarks and two production HPC codes (Quantum ESPRESSO, ICON dynamical core) showing up to $1.18\times$ speedup on GPUs and $1.12\times$ on CPUs.
\end{description}

\subsection{Computational Artifacts}

\begin{description}
\item[$A_1$] Concept DOI \url{https://doi.org/10.5281/zenodo.19827100} (resolves to the latest archived version). Two Git repositories back the artifact:
  \begin{itemize}
  \item \textbf{R1:} \url{https://github.com/ThrudPrimrose/SC26-Layout-AD}: microbenchmarks (E1--E5), ICON single-loopnest benchmarks (E6), and the full-module velocity-tendencies pipeline (E8).
  \item \textbf{R2:} \url{https://github.com/spcl/dace} (the DaCe framework; branch policy in §Software).
  \end{itemize}
\end{description}

\begin{center}
\small
\begin{tabular}{clll}
\toprule
Exp. & Contributions & Paper Elements & Repo \\
\midrule
E1 & $C_3$, $C_4$         & Fig.~4              & R1 \\
E2 & $C_{2\text{--}4}$    & Fig.~8              & R1 \\
E3 & $C_{2\text{--}4}$    & Fig.~9, Tab.~III    & R1 \\
E4 & $C_{2\text{--}4}$    & Fig.~10             & R1 \\
E5 & $C_{2\text{--}4}$    & Fig.~11, Lst.~1     & R1 \\
E6 & $C_3$, $C_4$         & Figs.~12--13, Tab.~IV & R1, R2 \\
E8 & $C_3$, $C_4$         & Fig.~14, Tab.~V     & R1, R2 \\
\bottomrule
\end{tabular}
\end{center}

\paragraph{E8 is the AD's default full velocity-tendencies path.}
\texttt{Experiments/E8\_LegacyVT/} reproduces Fig.~14 / Tab.~V from a single \texttt{sbatch run\_\{daint,beverin\}.sh}. The run script supports both a curated default sweep (paper figures, including E6's V$_k$ cross-product cells) and a raw \,cv\,$\times$\,ch\,$\times$\,f\,$\times$\,s\,$\times$\,n bit-product sweep; see $T_{8.3}$ for the config selection. \texttt{E7\_FullVelocityTendencies/} is an in-progress refactor and not on the AD path.

All figures regenerate via \texttt{Figures/plot\_all.sh}: quotient-graph figures (2--3) require no runtime input (illustrative); runtime figures (4, 8--14) read CSVs from \texttt{PaperSnapshot/} (frozen) or \texttt{Experiments/<experiment>/results/} (fresh) via \texttt{plot\_paper\_snapshot.sh} / \texttt{plot\_results.sh}. Hardware specs in §Artifact Setup.

%%%%%%%%%%%%%%%%%%%%%%%%%%%%%%%%%%%%%%%%%%%%%%%%%%%%%%%%
\section{Artifact Identification}
%%%%%%%%%%%%%%%%%%%%%%%%%%%%%%%%%%%%%%%%%%%%%%%%%%%%%%%%

\newartifact

\artrel

Artifact~$A_1$ covers all four contributions. The proof for $C_1$ ships as \texttt{Proofs/necessity\_proof.pdf} (read-only; not exercised by the software). The experiments validate the layout algebra ($C_2$), the cost metrics ($C_3$), and the cookbook-driven workflow ($C_4$): E1--E4 are microbenchmarks that isolate individual primitives, E5 evaluates composed transformations on a Quantum ESPRESSO production kernel, E6 applies the full optimization workflow per loop-nest pattern class, and E8 drives the resulting layouts through the velocity-tendencies module end-to-end on the GPU.

\artexp

Reproduction is qualitative. Reviewers should observe the same direction of effect per figure (which configuration wins, where the violin plots cluster); absolute bandwidths and gap magnitudes vary with hardware and system load.

\begin{itemize}\setlength{\itemsep}{1pt}\setlength{\parsep}{0pt}
\item \textbf{E1 (Fig.~4):} the permuted variant outperforms schedule-only across CPU and GPU; the GPU runs reach very close to the STREAM bandwidth peak.
\item \textbf{E2 (Fig.~8):} AoSoA-16 is competitive or best across backends. On CPUs, SoA at high $P$ over-saturates the cache-associativity limit and shows a performance collapse from page conflicts; AoS degrades more gradually. On GPUs, AoS, SoA, and AoSoA cluster more tightly.
\item \textbf{E3 (Fig.~9, Tab.~III):} the blocked layout beats naive row-major on CPUs and is competitive with state-of-the-art libraries. On GPUs the effect is muted: best-blocked still beats best-row-major, but the gap is much smaller.
\item \textbf{E4 (Fig.~10):} sorted indirection outperforms unsorted on every backend; the GPU gap is larger than the CPU gap.
\item \textbf{E5 (Fig.~11):} the composition unzip + permute + shuffle matches or beats the baseline on every backend.
\item \textbf{E6 (Figs.~12--13, Tab.~IV):} for each loopnest, at least one V$_k$ group (e.g.\ V6, the vertical-first variant) outperforms the ICON default; the indirect-stencil loopnest shows the most pronounced gain.
\item \textbf{E8 (Fig.~14, Tab.~V):} the vertical-first group (denoted V6) outperforms the connectivity-only (V2) and identity (V1) groups on Config~B.
\end{itemize}

\arttime
\begin{center}
\small
\begin{tabular}{lrr}
\toprule
Step & Daint & Beverin \\
\midrule
Setup                          &  15 &  15 \\
E1 Matrix Add.                 &  60 &  60 \\
E2 Conjugation                 &  60 &  60 \\
E3 Transpose                   &  60 &  60 \\
E4 GAS                         &  60 &  60 \\
E5 QE \texttt{addusxx\_g}      &  60 &  60 \\
E6 Loopnest 1                  & 120 & 120 \\
E6 Loopnests 2--6              & 600 & 600 \\
E8 Full-module sweep           & 1080 & 1440 \\
Analysis                       &  10 &  10 \\
\midrule
\textbf{Total}                 & \textbf{$\sim$35~hr} & \textbf{$\sim$41~hr} \\
\bottomrule
\end{tabular}
\end{center}
\noindent All times in minutes. E1--E5 are independent and can run in parallel.

\paragraph{Reproducibility caveat.} \%-of-STREAM annotations may differ slightly from the paper's. The per-platform STREAM Triad peak used for normalization is taken as the maximum of (a) E0\_NUMA's measured Triad and (b) the maximum sustained bandwidth observed across the E1--E4 binaries themselves; the cache-flush kernel was rewritten as a 3-sweep $8192^{2}$ Jacobi with buffer swap to defeat dead-code-elimination on GPUs, and the MI300A peak was re-measured against the corrected flush. Every figure in this artifact is normalized to that peak. Qualitative trends are unchanged: layout orderings and the gap between \texttt{unpermuted} and the V$_k$ winners reproduce.

\artin

\artinpart{Hardware}

Two HPC clusters at the Swiss National Supercomputing Centre (CSCS), accessed via SLURM with exclusive single-node allocations:

\begin{itemize}
\item \textbf{Daint.Alps:} Quad-GH200 nodes, 72-core Arm Neoverse~V2 Grace (512\,GB/s LPDDR5X) + NVIDIA H200 (4\,TB/s HBM3). Measured STREAM Triad peaks: CPU 1.61\,TB/s, GPU 3.78\,TB/s.
\item \textbf{Beverin:} Quad-MI300A APU nodes, 24 AMD Zen\,4 cores sharing 128\,GB unified HBM3. Measured STREAM Triad peaks: CPU 0.97\,TB/s, GPU 3.55\,TB/s.
\end{itemize}

\artinpart{Software}

Both platforms share Python (system \texttt{/usr/bin/python3.11}), DaCe, and the same canonical compile flags. Per-platform toolchain pins (Spack-managed):

\begin{center}
\small
\begin{tabular}{lll}
\toprule
Component & Daint.Alps & Beverin \\
\midrule
GCC          & 14.3              & 14.x              \\
GPU stack    & CUDA 12.9 / nvcc  & ROCm 6.4.1 / hipcc \\
GPU arch     & \texttt{sm\_90}   & \texttt{gfx942}   \\
OpenBLAS     & 0.3.29            & 0.3.30            \\
Tensor lib (E3) & cuTENSOR (spack) & hipTensor\footnotemark[1] (source) \\
Vendor xpose (E3) & HPTT\footnotemark[2] (source) & HPTT (source) \\
\bottomrule
\end{tabular}
\end{center}
\footnotetext[1]{hipTensor tag \texttt{rocm-6.4.1} (matching the ROCm pin): \url{https://github.com/ROCm/hipTensor}}
\footnotetext[2]{HPTT \texttt{v1.0.5}: \url{https://github.com/springer13/hptt}}

\noindent DaCe branch policy: E1--E6 do not import DaCe (pure C++/CUDA/HIP microbenchmarks plus plain-Python JSON crunching); E7 pins \texttt{yakup/dev}; E8 pins \texttt{f2dace/staging}. Each \texttt{run\_<pl>.sh} that needs DaCe exports \texttt{DACE\_BRANCH} explicitly before sourcing \texttt{activate.sh}; reviewers should not run E7 and E8 concurrently against the same DaCe clone.

\paragraph{CUDA pin.} The cluster is pinned to \texttt{cuda@12.9} via \texttt{Experiments/common/setup\_daint.sh}; 12.9 is the latest end-to-end-verified version. Experimental CUDA~13 support ships on \texttt{f2dace/staging} but is not part of the AD's verified path.

\paragraph{Compile flags and pinning.} \texttt{CPU\_CXXFLAGS} and \texttt{GPU\_CXXFLAGS} (exported by the setup scripts, applied uniformly across all \texttt{g++} / \texttt{nvcc} / \texttt{hipcc} invocations) carry \texttt{-O3 -march=native -ffast-math -fno-trapping-math -fno-math-errno -fopenmp -std=c++17} on both platforms (the no-trapping-math / no-math-errno flags are listed explicitly even though already implied by \texttt{-ffast-math}, and the equivalents are passed to nvcc via \texttt{-Xcompiler=\ldots} and to hipcc / clang natively); \texttt{-arch=sm\_90 --use\_fast\_math} on Daint; \texttt{--offload-arch=gfx942} with HIP-specific atomics on Beverin. \texttt{OMP\_SCHEDULE=static} and \texttt{SLURM\_CPU\_BIND=cores} are exported alongside; every \texttt{\#pragma omp parallel for} uses \texttt{schedule(static)}. Exact lists in each \texttt{setup\_<pl>.sh}.

\artinpart{Datasets / Inputs}

E1--E3 generate input data synthetically. Production-code experiments (E4, E5, E6, E8) auto-fetch their datasets via per-experiment \texttt{download\_data.sh} on first \texttt{sbatch}. URLs are public (ETH PolyBox, anonymous HTTP):

\begin{itemize}
\item \textbf{E4:} BaTiO$_3$ indirect-access patterns ($\sim$520\,KB) -- \url{https://polybox.ethz.ch/index.php/s/ySAEwfRTSH4Jzmr/download/bin_indirect_accesses-20260406T132129Z-1-001.zip}.
\item \textbf{E5:} BaTiO$_3$ from a QE \texttt{pw.x} HSE06 self-consistent calculation, 10-atom cell ($\sim$880\,MB) -- \url{https://polybox.ethz.ch/index.php/s/sBgGF5y4D2nbk25/download/bin_data_addusxx.tar.gz}.
\item \textbf{E6, E7, E8:} R02B05 ICON dataset, 81\,920 horizontal cells $\times$ 90 vertical levels ($\sim$5.8\,GB tarball, $\sim$9\,GB unpacked) -- \url{https://polybox.ethz.ch/index.php/s/WoxiditKJjgdEfR/download/nproma20480_data_files.tar.xz}. E6 and E8 each fetch independently; E7 can reuse a populated dir via \texttt{LOCAL\_DATA\_DIR}.
\end{itemize}

\artinpart{Installation and Deployment}

One-time per cluster: \texttt{bash Experiments/common/setup.sh} bootstraps a Python 3.11 venv (system \texttt{/usr/bin/python3.11}, with a \texttt{pyenv install 3.11.x} fallback if not present), the DaCe checkout, and the per-platform compile flags. Per experiment: \texttt{cd Experiments/<experiment> \&\& sbatch run\_<pl>.sh}, then \texttt{python plot\_paper.py} (or \texttt{bash Figures/plot\_all.sh Runtime} for the full sweep). Each run script sources \texttt{../common/activate.sh} (venv + DaCe pin) and \texttt{../common/setup\_<pl>.sh} (Spack loads, compile flags, OMP/NUMA pinning). \textbf{Throughout this AD, \texttt{<pl>} abbreviates the platform: \emph{daint} or \emph{beverin}.}

\artcomp

Every experiment is driven by its \texttt{run\_<pl>.sh}, which compiles the benchmark, runs 100 timed repetitions after 5 warm-ups, and writes results next to the experiment, with no intermediate stage invoked by hand. Plotting is invoked manually with \texttt{python plot\_paper.py} once the job completes. E1--E5 are mutually independent and may be submitted in parallel.

\paragraph{Per-experiment workflow.}
\begin{itemize}\setlength{\itemsep}{1pt}\setlength{\parsep}{0pt}
\item \textbf{E1} (Fig.~4): \texttt{E1\_MatrixAdd/run\_<pl>.sh} compiles \texttt{bench\_\{cpu,gpu\}}, runs the schedule-only vs.\ permuted sweep; \texttt{plot\_paper.py} emits Fig.~4.
\item \textbf{E2} (Fig.~8): \texttt{E2\_Conjugation/run\_<pl>.sh} sweeps $P$ over AoS / SoA / AoSoA-16; \texttt{plot\_paper.py} emits Fig.~8.
\item \textbf{E3} (Fig.~9, Tab.~III): \texttt{E3\_Transpose/run\_<pl>.sh} runs the transpose sweep against HPTT / cuTENSOR / hipTensor and a follow-up \texttt{cost\_metrics} pass for $\mu$ and $\Delta$; \texttt{plot\_paper.py} emits Fig.~9 and Tab.~III.
\item \textbf{E4} (Fig.~10): \texttt{E4\_GAS/run\_<pl>.sh} auto-fetches BaTiO$_3$ if absent, runs uniform-random and index-sort gather--accumulate--scatter; \texttt{plot\_paper.py} emits Fig.~10.
\item \textbf{E5} (Fig.~11): \texttt{E5\_USXX/run\_<pl>.sh} auto-fetches BaTiO$_3$ HSE06 state if absent, runs the composed unzip + permute + shuffle on \texttt{addusxx\_g}; \texttt{plot\_paper.py} emits Fig.~11.
\item \textbf{E6} (Figs.~12--13, Tab.~IV): each \texttt{loopnest\_\{1..6\}/run\_<pl>.sh} (under \texttt{E6\_VelocityTendencies/}) automates the full per-loopnest pipeline:
  \begin{itemize}\itemsep1pt
  \item $T_{6.1}$ data fetch (R02B05 ICON state).
  \item $T_{6.2}$ LOKI access analysis $\to$ \texttt{layout\_candidates.json}.
  \item $T_{6.3}$ 6-group taxonomy (\textit{cv, ch, f, s, n, lm}) $\to$ \texttt{canonical\_array\_groups.json}.
  \item $T_{6.4}$ layout sweep over $nlev \in \{90, 96, 128, 256\}$, default 128.
  \item $T_{6.5}$ winners cross-product (\texttt{generate\_winners.py} / \texttt{derive\_winners.py} / \texttt{generate\_layout\_groups.py}) $\to$ \texttt{winners.json}, \texttt{layout\_groups.json}.
  \item $T_{6.6}$ \texttt{plot\_paper.py} $\to$ Figs.~12--13, Tab.~IV.
  \end{itemize}
\item \textbf{E8} (Fig.~14, Tab.~V): \texttt{E8\_LegacyVT/run\_<pl>.sh} consumes \texttt{winners.json} (from E6's $T_{6.5}$, or auto-regenerated on the fly if missing) and automates:
  \begin{itemize}\itemsep1pt
  \item $T_{8.1}$ dataset resolution (existing E8 dir $\to$ E7 symlink $\to$ fresh fetch).
  \item $T_{8.2}$ per-config compile via \texttt{compile\_gpu\_stage8}.
  \item $T_{8.3}$ run \texttt{run\_stage8\_permutations.py} for each \texttt{CONFIGS} entry, in one of three modes:
    \begin{itemize}\setlength{\itemsep}{0pt}\setlength{\parsep}{0pt}
    \item \emph{curated} (default): 67 cells / 134 binaries, namely \texttt{nlev\_first}, \texttt{index\_only}, \texttt{winner\_v1}, plus the 64 unique \texttt{v123\_*} cross-product cells from E6's \texttt{layout\_crossproduct\_v123.json}.
    \item \emph{auto-generated}: the full $\,$cv\,$\times$\,ch\,$\times$\,f\,$\times$\,s\,$\times$\,n raw sweep, via \texttt{CONFIGS=all}.
    \end{itemize}
    Stdout per (config $\times$ shuffle $\times$ timestep) lands in \texttt{<pl>\_full\_permutations\_8/} as one TXT each.
  \item $T_{8.4}$ plot via \texttt{plot\_paper\_snapshot.sh} or \texttt{plot\_results.sh} (under \texttt{Figures/}).
  \end{itemize}
\end{itemize}

\paragraph{Inspecting configs before sbatch.} Each run script accepts a \texttt{--dry-run} flag that lists every config that would be compiled and run without touching the GPU; E7 exposes the same via \texttt{python tools/run\_layout\_configs.py --dry-run}. The named \texttt{winner\_v\{1,2,6\}} aliases are byte-identical on both paths and consume \texttt{layout\_groups.json} from E6 at the canonical \emph{nlev=128} slice.

\paragraph{Statistical methodology.} Following Hoefler and Belli, every configuration runs 100 repetitions after 5 warm-ups, with cache flushes between repetitions (an $8192^2$ Jacobi stencil over 3 sweeps with buffer swap; \texttt{common/jacobi\_flush.h}). We report medians with 95\% BCa-bootstrap confidence intervals ($n=10{,}000$ resamples). Violin plots show all 100 samples without outlier trimming; every violin uses the same KDE settings (Scott's rule, 200 evaluation points). The RNG seed is \texttt{SC26\_SEED=42}, and bandwidths are normalized to the measured STREAM Triad peak per platform.

\artout

Each experiment writes results into its own folder: \texttt{results/\{daint,beverin\}/} for the CSV-emitting binaries (E1--E6), and \texttt{<pl>\_full\_permutations\_8/} for E8's per-config stage-8 TXTs. \texttt{plot\_paper.py} then produces the PDF violin plots matching Figures~4 and 8--14; the metric tables (Tabs.~III, IV, V) are emitted alongside. Expected trends reproduce: relative layout ordering and gap magnitudes match the paper's. Absolute bandwidths depend on hardware and system load.

\end{document}
